\documentclass[12pt]{article}
\usepackage[margin=1in]{geometry}
\usepackage{graphicx}
\usepackage{booktabs}
\usepackage{amsmath}

\title{The Effect of Paid Search Advertising on Revenue: \\
A Replication of the eBay Natural Experiment}
\author{Caroline Allman}
\date{\today}

\begin{document}
\maketitle

\section{Introduction}
This paper asks: what is the causal effect of paid search advertising on eBay's revenue? Firms spend
large amounts on search engine marketing (SEM), but measuring its true impact is difficult because
advertising is typically targeted toward markets or times when demand is already high. To address this
endogeneity problem, we rely on a natural experiment described by Blake, Nosko, and Tadelis (2014),
in which eBay temporarily turned off paid search ads in selected markets. Using a difference-in-differences
(DID) design, we estimate how revenue changed in treated markets relative to untreated markets during
the advertising shutdown.

\section{Experimental Setting}
Blake et al. (2014) worked with eBay to pause bidding on Google AdWords in a subset of U.S. designated
market areas (DMAs). Treatment began on May 22, 2012 and lasted for eight weeks. In treated DMAs,
paid search was turned off (search\_stays\_on = 0), while in control DMAs it remained on
(search\_stays\_on = 1). The experiment was not fully randomized: the largest DMAs were excluded from
treatment, which creates systematic baseline differences in revenue between the treatment and control groups.

Because treated and control DMAs differ in levels before the intervention, a simple comparison of average
revenue is not appropriate. Instead, we use a DID approach that compares changes over time in treated
markets to changes over time in control markets. Under the parallel trends assumption, the DID estimator
isolates the effect of turning off paid search advertising.

\section{Data and Figures}
The dataset contains daily revenue observations for 210 DMAs from April through July 2012. Key variables
include the date, DMA identifier, treatment group, post-treatment period indicator, and revenue. The
analysis first visualizes average revenue patterns for the treatment and control groups and then estimates
a DID model.

\begin{figure}[h]
\centering
\includegraphics[width=0.85\textwidth]{../output/figures/figure_5_2.png}
\caption{Average revenue for treatment and control DMAs over time.
The dashed line marks May 22, 2012, when SEM was turned off
for the treatment group.}
\label{fig:revenue}
\end{figure}

Figure \ref{fig:revenue} shows that the control group has higher average revenue throughout the sample,
which is consistent with the design choice to exclude the largest DMAs from treatment. Both groups exhibit
similar week-to-week oscillation patterns, suggesting shared seasonality and demand cycles. At this scale,
there is no dramatic divergence visible at the treatment date, which motivates looking at a transformed
version of the series and then estimating the DID formally.

\begin{figure}[h]
\centering
\includegraphics[width=0.85\textwidth]{../output/figures/figure_5_3.png}
\caption{Log-scale average revenue difference between control and treatment
groups over time. The gap appears to increase after May 22.}
\label{fig:logdiff}
\end{figure}

Figure \ref{fig:logdiff} plots the log-scale difference (equivalently, the log ratio) between control and
treatment group average revenues. Before May 22, the series fluctuates around a relatively stable level.
After the intervention, the gap appears to increase slightly, which is consistent with treated markets
experiencing a modest decline in revenue relative to controls. This visual pattern provides motivation for
the DID regression results reported next.

\section{Results}
\input{../output/tables/did_table.tex}

The DID estimate in logs is negative and small: $\hat{\gamma} \approx -0.0066$. Interpreting this in levels,
$\exp(\hat{\gamma}) \approx 0.9934$, which corresponds to about a 0.66\% decline in revenue in treatment
DMAs when paid search was turned off. However, the 95\% confidence interval includes zero
(approximately $[-0.0175, 0.0043]$), implying the effect is not statistically significant at the 5\% level.
Economically, this suggests that paid search has little measurable effect on revenue for a well-known brand
like eBay, potentially because many users would find eBay through organic search even without paid ads.

\section{Conclusion}
Using a difference-in-differences design and eBay's natural experiment, we find a small and statistically
insignificant reduction in revenue when paid search advertising is turned off. This replicates the core
finding of Blake et al. (2014): for an established brand, the incremental return to paid search is limited.
Important limitations are that the revenue data are scaled/shifted rather than raw dollar amounts, the
treatment assignment was not fully randomized, and the analysis uses a simple DID framework rather than
a richer regression specification. Still, the results highlight why causal identification is critical when
evaluating marketing ROI.

\begin{thebibliography}{9}
\bibitem{blake2014}
Blake, T., C. Nosko, and S. Tadelis (2014).
Consumer Heterogeneity and Paid Search Effectiveness: A Large-Scale Field Experiment.
\textit{Econometrica}, 83(1): 155--174.

\bibitem{taddy2019}
Taddy, M. (2019).
\textit{Business Data Science}. McGraw-Hill, Chapter 5.
\end{thebibliography}

\end{document}
